\documentclass[
	% -- opções da classe memoir --
	12pt,				% tamanho da fonte
	openright,			% capítulos começam em pág ímpar (insere página vazia caso preciso)
	oneside,			% para impressão em um lado. Oposto a twoside, verso e anverso
	a4paper,			% tamanho do papel.
	% -- opções da classe abntex2 --
	chapter=TITLE,		% títulos de capítulos convertidos em letras maiúsculas
	section=TITLE,		% títulos de seções convertidos em letras maiúsculas
	%subsection=TITLE,	% títulos de subseções convertidos em letras maiúsculas
	%subsubsection=TITLE,% títulos de subsubseções convertidos em letras maiúsculas
	% -- opções do pacote babel --
	english,			% idioma adicional para hifenização
	%french,				% idioma adicional para hifenização
	%spanish,			% idioma adicional para hifenização
	brazil,				% o último idioma é o principal do documento
	]{styles/abntex2}

% \usepackage[alf]{styles/abntex2cite}	% Citações padrão ABNT

\usepackage[alf,
abnt-emphasize = bf, % destaca o titulo da revista ou livro em negrito;
abnt-etal-list = 0, % trabalhos com mais de 3 autores recebem et al.,;
%abnt-etal-text = it, % escreve o et al., em italico;
abnt-and-type = e, % usa o carater '&' no lugar de 'e' para mais de um autor;
abnt-last-names = abnt, % trata sobrenomes 'estritamente' conforme a ABNT; e
abnt-repeated-author-omit = no, % autores com + de uma entrada recebem '__.'
%bibjustify
]{styles/abntex2cite}


\usepackage[utf8]{inputenc}
\usepackage{listings}
\usepackage{lastpage}
\usepackage{times}
\usepackage{graphicx}

\usepackage{tabularx}
\usepackage{makecell}
\usepackage{lmodern}
\usepackage{caption}
\usepackage{subcaption}
\captionsetup{justification   = raggedright,
              singlelinecheck = false}
\usepackage[table,xcdraw]{xcolor}
\usepackage[normalem]{ulem}
\usepackage{amsmath}


% \usepackage{natbib}

\usepackage{xcolor}
\usepackage{enumerate}

\usepackage{styles/ttc_univali}

\title{Trabalho Técnico Científico}
\author{Gabriel Sodré da Silva Prates}
\date{Novembro, 2026}

\sloppy

\newcommand{\citet}{\citeonline}

\usepackage{pgfplots}

% Estilo para trechos de código TypeScript (frontend e backend Node).
\lstdefinestyle{ts-style}{
  language=Java, % usa highlight de Java como fallback (mais próximo de TS)
  basicstyle=\ttfamily\small,
  keywordstyle=\color{blue},
  stringstyle=\color{orange},
  commentstyle=\color{gray},
  backgroundcolor=\color{white},
  frame=single,
  breaklines=true,
  showstringspaces=false
}

\lstdefinestyle{python-style}{
  language=Python,
  basicstyle=\ttfamily\small,
  keywordstyle=\color{blue},
  stringstyle=\color{orange},
  commentstyle=\color{gray},
  backgroundcolor=\color{white},
  frame=single,
  breaklines=true,
  showstringspaces=false
}

\begin{document}

\pretextual

\begin{Info}
% Universidade
{UNIVERSIDADE DO VALE DO ITAJAÍ}
% Escola
{ESCOLA POLITÉCNICA}
% Curso
{CURSO DE CIÊNCIA DA COMPUTAÇÃO}
% Titulo
% TODO: validar com o orientador. Alternativas em discussao:
%   (A) ...COM EXTRACAO VIA LLM
%   (B) ...COM EXTRACAO ESTRUTURADA VIA LLM
%   (C) ...COM PIPELINE SERVERLESS E EXTRACAO VIA LLM
{ROLÊ SC: APLICATIVO MOBILE AGREGADOR DE EVENTOS CULTURAIS E MUNICIPAIS DO LITORAL NORTE DE SANTA CATARINA COM EXTRAÇÃO AUTOMATIZADA VIA LLM}
% Autor
{Gabriel Sodré da Silva Prates}
% Cidade e Data
{Itajaí (SC), novembro de 2026}
% Nome da Área de concentração
{Engenharia de Software}
% Orientador(a)
{Alex Luciano Roesler Rese, Me.}
{}
\end{Info}


\begin{Resumo}
A descoberta de eventos culturais, esportivos, gastronômicos e comunitários no litoral norte de Santa Catarina enfrenta um problema de fragmentação: cada prefeitura, fundação cultural e produtor publica suas programações em sites próprios e redes sociais, sem padronização nem agregação regional. Este trabalho de conclusão de curso (TCC) apresenta o desenvolvimento do \textit{Rolê SC}, um aplicativo móvel agregador de eventos dos municípios de Itajaí, Camboriú, Balneário Camboriú, Porto Belo, Itapema e Bombinhas, sustentado por um \textit{pipeline serverless} que automatiza a coleta e a estruturação dessas informações. O ecossistema é composto por quatro componentes independentes implantados na nuvem da \textit{Amazon Web Services}: um \textit{web crawler} baseado em \textit{Playwright} que coleta páginas dos sites oficiais; um processador de linguagem natural que utiliza o modelo \textit{Llama 3.1 8B Instruct}, servido pela plataforma \textit{Groq}, para extrair eventos estruturados ou descartar páginas que não correspondem a eventos; uma \textit{API REST} em \textit{TypeScript} sobre o \textit{framework Fastify}, autenticada por \textit{JWT}; e um aplicativo móvel Android construído em \textit{React Native} com \textit{Expo}. O trabalho contempla também o estudo do modelo de negócio, a análise de concorrentes, o levantamento de requisitos e a documentação do sistema. Nesta primeira etapa (TCC-II), foram implementados o \textit{pipeline} de coleta e extração para os cinco municípios-alvo, os fluxos de cadastro, verificação de email, autenticação, recuperação e alteração de senha, edição de perfil e \textit{upload} de avatar, bem como a publicação do aplicativo Android via \textit{Expo Application Services}. Para o TCC-III, está planejada a integração ponta a ponta da autenticação federada com \textit{Google}, a implementação das funcionalidades sociais (curtidas, indicação de presença e comentários) e a avaliação empírica de desempenho, custo e qualidade da extração estruturada.

\textbf{Palavras-chave}: Aplicativo Móvel. Computação \textit{Serverless}. Extração de Informação. Modelos de Linguagem. \textit{Web Crawling}.
\end{Resumo}

\begin{Abstract}
The discovery of cultural, sports, gastronomic, and community events in the northern coast of Santa Catarina faces a chronic fragmentation problem: each municipality, cultural foundation, and independent producer publishes its programming on separate websites, social networks, and printed materials, without structural standardization or regional aggregation. This undergraduate thesis presents the development of \textit{Rolê SC}, a mobile event aggregator application for the municipalities of Itajaí, Camboriú, Balneário Camboriú, Porto Belo, Itapema, and Bombinhas, supported by a serverless pipeline that automates the collection and structuring of this information. The ecosystem consists of four independent components deployed on \textit{Amazon Web Services}: a \textit{web crawler} based on \textit{Playwright} that collects pages from official websites; a natural language processor that uses the \textit{Llama 3.1 8B Instruct} model, served by the \textit{Groq} platform, to extract structured events or to discard pages that do not correspond to events; a \textit{REST API} developed in \textit{TypeScript} on the \textit{Fastify} framework, authenticated by \textit{JWT}; and an Android mobile application built with \textit{React Native} and \textit{Expo}. The work also covers the business model study, the analysis of direct and indirect competitors, requirements gathering, and system documentation. In this first stage (TCC-II), the collection and extraction pipeline was implemented for the five target municipalities, along with the flows for registration, email verification, authentication, password recovery and change, profile editing, and avatar upload, as well as the publication of the Android application through \textit{Expo Application Services}. For TCC-III, the end-to-end integration of federated authentication with \textit{Google}, the implementation of social features (likes, attendance indication, and comments), and the empirical evaluation of performance, cost, and structured extraction quality are planned.

\textbf{Keywords}: Mobile Application. \textit{Serverless} Computing. Information Extraction. Language Models. \textit{Web Crawling}.
\end{Abstract}


\pdfbookmark[0]{\listfigurename}{lof}
\listoffigures*
\cleardoublepage

\pdfbookmark[0]{\listtablename}{lot}
\listoftables*
\cleardoublepage


%\pdfbookmark[0]{\listofequacaoname}{loe}
%\listofequacao*
%\cleardoublepage

% --- Lista de siglas e abreviaturas ---
\begin{siglas}
  \item[ABNT] Associação Brasileira de Normas Técnicas
  \item[API] \textit{Application Programming Interface}
  \item[APK] \textit{Android Package}
  \item[AWS] \textit{Amazon Web Services}
  \item[CDN] \textit{Content Delivery Network}
  \item[CORS] \textit{Cross-Origin Resource Sharing}
  \item[CRUD] \textit{Create, Read, Update, Delete}
  \item[DOM] \textit{Document Object Model}
  \item[EAS] \textit{Expo Application Services}
  \item[HTML] \textit{HyperText Markup Language}
  \item[HTTP] \textit{HyperText Transfer Protocol}
  \item[HTTPS] \textit{HyperText Transfer Protocol Secure}
  \item[IAM] \textit{Identity and Access Management}
  \item[IaC] \textit{Infrastructure as Code}
  \item[JSON] \textit{JavaScript Object Notation}
  \item[JWT] \textit{JSON Web Token}
  \item[LGPD] Lei Geral de Proteção de Dados
  \item[LLM] \textit{Large Language Model}
  \item[MVP] \textit{Minimum Viable Product}
  \item[NLP] \textit{Natural Language Processing}
  \item[NoSQL] \textit{Not only SQL}
  \item[OAuth] \textit{Open Authorization}
  \item[OIDC] \textit{OpenID Connect}
  \item[OTA] \textit{Over-the-Air}
  \item[RF] Requisito Funcional
  \item[RN] Regra de Negócio
  \item[RNF] Requisito Não Funcional
  \item[REST] \textit{Representational State Transfer}
  \item[S3] \textit{Simple Storage Service}
  \item[SAM] \textit{Serverless Application Model}
  \item[SC] Santa Catarina
  \item[SDK] \textit{Software Development Kit}
  \item[SES] \textit{Simple Email Service}
  \item[SLA] \textit{Service Level Agreement}
  \item[SWOT] \textit{Strengths, Weaknesses, Opportunities, Threats}
  \item[TCC] Trabalho de Conclusão de Curso
  \item[TLS] \textit{Transport Layer Security}
  \item[TTL] \textit{Time to Live}
  \item[UC] Caso de Uso
  \item[UML] \textit{Unified Modeling Language}
  \item[URL] \textit{Uniform Resource Locator}
  \item[UX] \textit{User Experience}
\end{siglas}


% ---
\pdfbookmark[0]{\contentsname}{toc}
\tableofcontents*
\cleardoublepage
% ---

 \clearpage


\textual % Indica inicio dos elementos textuais
\chapter{Modelo de negócios}
% =========================================================================
% Capítulo 1 — MODELO DE NEGÓCIOS
% =========================================================================
% Estrutura esperada (espelho do template UNIVALI):
%   1.1 Introdução ao domínio de aplicação
%   1.2 Público alvo
%   1.3 Solução Proposta (com Canvas e SWOT como figuras em images/)
%   1.4 Delimitação do Escopo
%   1.5 Levantamento de produtos existentes (concorrentes)
%       1.5.1 Sympla
%       1.5.2 Eventbrite
%       1.5.3 Facebook Events
%       1.5.4 Instagram (descoberta orgânica)
%       1.5.5 Apps municipais ("Visite Bombinhas", etc.)
%   1.6 Objetivos
%       1.6.1 Objetivos Gerais
%       1.6.2 Objetivos Específicos
%
% Base de informação:
%   - Seção 1, 2, 3, 4, 5 do docs/CONTEXTO_TCC_MONOGRAFIA.md
%   - Imagens necessárias em images/:
%       - ModeloNegocio.pdf  (Business Model Canvas)
%       - AnaliseSWOT.pdf
% =========================================================================

\section{Introdução ao domínio de aplicação}
% TODO: solicitar ao Claude com base na seção 2 do contexto.

\section{Público alvo}
% TODO: solicitar com base na seção 4 (Personas) do contexto.

\section{Solução Proposta}
% TODO: incluir aqui as figuras do Canvas e SWOT (precisam ser desenhadas em
% draw.io / Miro / Canva e exportadas para PDF em images/).

\section{Delimitação do Escopo}
% TODO: delimitar o que o MVP entrega vs. o que fica para o TCC III.
% Base: seção 15 (matriz de status) do contexto.

\section{Levantamento de produtos existentes}
% TODO: subseções para cada concorrente da seção 5 do contexto:
% Sympla, Eventbrite, Facebook Events, Instagram, apps municipais.

\section{Objetivos}
\subsection{Objetivos Gerais}
% TODO

\subsection{Objetivos Específicos}
% TODO: itemize com 5-7 objetivos específicos.


\chapter{Especificação do Produto}
% =========================================================================
% Capítulo 2 — ESPECIFICAÇÃO DO PRODUTO
% =========================================================================
% Estrutura esperada:
%   2.1 Tecnologias utilizadas
%       2.1.1 React Native + Expo (frontend mobile)
%       2.1.2 TypeScript
%       2.1.3 Fastify (backend API)
%       2.1.4 AWS Lambda e API Gateway (infraestrutura serverless)
%       2.1.5 MongoDB Atlas (persistência)
%       2.1.6 JWT e OAuth 2.0 / OIDC (autenticação)
%       2.1.7 Playwright (web crawling)
%       2.1.8 LLMs e Groq Cloud (Llama 3.1 8B para extração de eventos)
%       2.1.9 AWS S3 e SES (storage de avatares e envio de emails)
%   2.2 Documentação do sistema
%       2.2.1 Regras de Negócio (RN01-RN25) — usar \begin{quadro}
%       2.2.2 Requisitos Funcionais (RF01-RF32) — usar \begin{table}
%       2.2.3 Requisitos Não Funcionais (RNF01-RNF26) — usar \begin{table}
%       2.2.4 Casos de Uso — diagrama + quadros UC-VISITANTE-XX, UC-USUARIO-XX
%   2.3 Modelagem do banco de dados (diagrama ER + descrição de coleções)
%
% Base:
%   - Seção 6 (tecnologias justificadas), 8 (banco), 9 (RF), 10 (RNF),
%     11 (RN), 12 (casos de uso) do docs/CONTEXTO_TCC_MONOGRAFIA.md
%
% Imagens necessárias em images/:
%   - DiagramaCasosUsoVisitante.pdf
%   - DiagramaCasosUsoUsuario.pdf
%   - DiagramaER.pdf
% =========================================================================

\section{Tecnologias utilizadas}
% TODO: uma subseção por tecnologia, no padrão:
% (i) descrição da tecnologia (com citação),
% (ii) características-chave,
% (iii) papel no projeto Rolê SC e justificativa de escolha.

\section{Documentação do sistema}

\subsection{Regras de Negócio}
% TODO: \begin{quadro} com RN01-RN25 (seção 11 do contexto).

\subsection{Requisitos Funcionais}
% TODO: \begin{table} com RF01-RF32 (seção 9 do contexto), separadas por
% domínio (eventos, autenticação, perfil, favoritos).

\subsection{Requisitos não funcionais}
% TODO: \begin{table} com RNF01-RNF26 (seção 10 do contexto).

\subsection{Casos de Uso}
% TODO: diagrama + quadros detalhados dos UCs principais (UC01, UC04, UC07,
% UC10, UC11, UC14) da seção 12 do contexto.

\section{Modelagem do banco de dados}
% TODO: diagrama ER + descrição das coleções (seção 8 do contexto).


\chapter{Desenvolvimento Parcial}
% =========================================================================
% Capítulo 3 — DESENVOLVIMENTO PARCIAL
% =========================================================================
% Estrutura esperada:
%   3.1 Visão geral da arquitetura
%       (Diagrama de arquitetura geral + diagramas de sequência:
%        login Google e recuperação de senha)
%   3.2 Pipeline de coleta (role-sc-web-crawler)
%   3.3 Processador NLP (role-sc-processador-nlp-webcrawls)
%   3.4 Backend REST (role-sc-api)
%       3.4.1 Endpoints e contrato
%       3.4.2 Autenticação JWT
%       3.4.3 Integração com S3 (avatar) e SES (emails)
%   3.5 Aplicativo Mobile (role-sc-frontend)
%       3.5.1 Arquitetura (AuthContext, Navigation, Services)
%       3.5.2 Telas (Login, Cadastro, Verificar Email, Recuperar Senha,
%             Confirmar Código, Redefinir Senha, Home, Detalhe Evento,
%             Favoritos, Perfil, Editar Perfil, Alterar Senha)
%       3.5.3 Integração OAuth com Google
%       3.5.4 Upload de avatar via presigned URL
%   3.6 Publicação (EAS Build, APK preview, planejamento Play Store)
%
% Base:
%   - Seções 7 (arquitetura), 13 (protótipos das telas), 6 (tecnologias) do
%     docs/CONTEXTO_TCC_MONOGRAFIA.md
%   - Todos os docs/CONTRATO_*.md para detalhes técnicos por fluxo
%
% Imagens necessárias em images/ (capturas reais do APK preview):
%   - ArquiteturaGeral.pdf            (diagrama do Mermaid exportado)
%   - SequenciaLoginGoogle.pdf
%   - SequenciaRecuperacaoSenha.pdf
%   - TelaLogin.png
%   - TelaCadastro.png
%   - TelaVerificarEmail.png
%   - TelaConfirmarCodigo.png
%   - TelaHome.png
%   - TelaHomeFiltrada.png
%   - TelaDetalheEvento.png
%   - TelaFavoritos.png
%   - TelaPerfil.png
%   - TelaEditarPerfil.png
%   - TelaAlterarSenha.png
% =========================================================================

\section{Visão geral da arquitetura}
% TODO

\section{Pipeline de coleta}
% TODO

\section{Processador NLP}
% TODO

\section{Backend REST}
% TODO

\section{Aplicativo Mobile}
% TODO: subseções com cada tela, descrição funcional + captura

\section{Publicação}
% TODO: descrever EAS, fluxo de build APK/AAB, planejamento Play Store
% (base: docs/PUBLICACAO.md)


\chapter{Resultados parciais}
% =========================================================================
% Capítulo 4 — RESULTADOS PARCIAIS
% =========================================================================
% Estrutura esperada (espelhando o Thiago):
%   4.1 Cenário de teste (ambientes, dispositivos, dados de entrada)
%   4.2 Testes Funcionais
%       4.2.1 Validação do pipeline (crawler -> NLP -> eventos)
%       4.2.2 Testes da API REST
%       4.2.3 Testes do aplicativo mobile (telas e fluxos)
%       4.2.4 Teste final do MVP (cenário fim-a-fim)
%   4.3 Métricas observadas (quantidade de eventos coletados, latência,
%       precisão da extração da IA, taxa de sucesso de login, etc.)
%
% Base:
%   - Seção 15 (status de implementação) do CONTEXTO_TCC_MONOGRAFIA.md
%   - Você (Gabriel) precisa MEDIR e capturar dados reais para essa seção:
%       (a) quantos eventos foram coletados por município ao longo de N dias
%       (b) acurácia da classificação isEvent=true/false (amostragem manual)
%       (c) latência média de GET /eventos (p50, p95, p99)
%       (d) tempo médio de cold start vs warm start da Lambda
%       (e) cobertura de testes manuais por tela
% =========================================================================

\section{Cenário de teste}
% TODO

\section{Testes Funcionais}
% TODO: subseções por fase de validação

\section{Métricas observadas}
% TODO: tabelas com números reais coletados


\chapter{Análise crítica}
% =========================================================================
% Capítulo 5 — ANÁLISE CRÍTICA
% =========================================================================
% Estrutura esperada (espelhando o Thiago):
%   5.1 Análise dos resultados (discussão crítica do que funcionou e não
%       funcionou no MVP — pontos de melhoria, limitações observadas)
%   5.2 Planejamento do TCC III
%       5.2.1 Objetivos para o TCC III
%       5.2.2 Cronograma
%       5.2.3 Metodologia
%       5.2.4 Análise de riscos
%
% Base:
%   - Seção 15 (status) e 16 (roadmap) do CONTEXTO_TCC_MONOGRAFIA.md
%   - Itens para o TCC III, partindo do roadmap:
%       * Adicionar crawlers de Porto Belo, Itapema e Bombinhas
%       * Implementar notificações push
%       * Implementar curtidas, comentários e "eu vou"
%       * Adicionar métricas de uso e analytics
%       * Publicar na Play Store (track interno)
%       * Conformidade LGPD (política de privacidade, exclusão de conta)
%       * Implementar token versioning (invalidação de JWT na troca de senha)
% =========================================================================

\section{Análise dos resultados}
% TODO

\section{Planejamento do TCC III}

\subsection{Cronograma}
% TODO: tabela mês a mês (fev/2027 a maio/2027) com as atividades
% (modelo do Thiago, resultados.tex linhas 58-80)

\subsection{Metodologia}
% TODO

\subsection{Análise de riscos}
% TODO: tabela com Risco / Probabilidade / Impacto / Gatilho / Contingência
% (modelo do Thiago, resultados.tex linhas 92-120)


\pagestyle{simple} % remove o cabecalho



\postextual

\bibliography{referencias}

% Glossário (pós-textual, depois das referências — ABNT NBR 14724/2011)
% Glossário do TCC do Rolê SC.
% Posicionamento no main.tex: depois de \bibliography{referencias},
% dentro do bloco \postextual (ordem ABNT NBR 14724/2011).
% Comando para inclusão: % Glossário do TCC do Rolê SC.
% Posicionamento no main.tex: depois de \bibliography{referencias},
% dentro do bloco \postextual (ordem ABNT NBR 14724/2011).
% Comando para inclusão: % Glossário do TCC do Rolê SC.
% Posicionamento no main.tex: depois de \bibliography{referencias},
% dentro do bloco \postextual (ordem ABNT NBR 14724/2011).
% Comando para inclusão: \input{glossario.tex}

\pretextualchapter{Glossário}

\begin{description}

\item[\textit{Antienumeration}] Padrão de segurança em que respostas a entradas inexistentes (\textit{e.g.}, \textit{e-mail} não cadastrado) são mantidas idênticas às de entradas válidas, evitando que um atacante descubra usuários cadastrados por meio de respostas distintas.

\item[\textit{API Key}] Cabeçalho HTTP (\texttt{x-api-key}) validado pelo \textit{API Gateway} antes da invocação da \textit{Lambda}, utilizado para aplicar \textit{throttling} e quota.

\item[\textit{AsyncStorage}] \textit{API} de armazenamento \textit{key-value} persistente do \textit{React Native}, utilizada no aplicativo para guardar o \textit{JWT}, o usuário autenticado e a lista de favoritos.

\item[\textit{Cold start}] Tempo adicional de inicialização sofrido por uma \textit{Lambda} quando não há contêiner aquecido disponível, antes do código de usuário começar a executar.

\item[Evento] Atividade pública (cultural, esportiva, gastronômica, religiosa ou comunitária) realizada em um dos seis municípios-alvo do projeto, identificada e estruturada pelo \textit{pipeline} de extração via \textit{LLM}.

\item[\textit{id\_token}] Credencial no formato \textit{JWT}, assinada em \textit{RS256}, emitida pelo \textit{Google} após autenticação \textit{OAuth} bem-sucedida; o \textit{backend} valida sua assinatura e troca-a pelo \textit{JWT} do próprio sistema.

\item[\textit{isEvent}] Atributo booleano definido pelo processador de linguagem natural para cada documento bruto: \texttt{true} indica que a página foi classificada como um evento e fica visível no aplicativo; \texttt{false} indica notícia comum, descartada pela \textit{API}.

\item[\textit{JWT}] Credencial \textit{stateless} no formato \textit{JSON Web Token}, assinada em \textit{HS256}, que prova a identidade do usuário ao \textit{backend} por sete dias após o \textit{login}.

\item[\textit{Linking} automático] Comportamento do \textit{login} com \textit{Google} em que, se o \textit{e-mail} do \textit{id\_token} já corresponde a uma conta existente, o identificador \textit{googleId} é vinculado a essa conta sem solicitar senha.

\item[\textit{Pipeline}] Sequência de componentes \textit{serverless} do Rolê SC que executa, de forma encadeada e idempotente, a coleta dos sites oficiais, o armazenamento bruto e a extração estruturada de eventos.

\item[\textit{Presigned URL}] \textit{URL} temporária do \textit{Amazon S3} (com tempo de vida de cinco minutos) assinada pelo \textit{backend}, que autoriza o aplicativo a executar um \textit{PUT} diretamente no \textit{bucket} sem que o conteúdo do arquivo passe pelo servidor.

\item[\textit{Raw crawl}] Documento bruto (\textit{HTML} ou texto) capturado de um \textit{site} público pelo \textit{web crawler} e armazenado na coleção \texttt{raw\_crawls} antes do processamento de linguagem natural.

\item[\textit{Serverless}] Modelo de execução em nuvem no qual o provedor gerencia integralmente a infraestrutura subjacente, e o consumidor é cobrado apenas pelo tempo efetivo de execução das funções, sem provisionar servidores.

\item[\textit{Strength meter}] Componente visual que apresenta ao usuário, em tempo real, quais requisitos de força da senha (comprimento, presença de caracteres especiais, dígitos etc.) já foram atendidos.

\item[\textit{TTL Index}] Índice especial do \textit{MongoDB} que remove automaticamente, em segundo plano, documentos cujo campo de expiração já passou, utilizado neste trabalho para tokens efêmeros de verificação de \textit{e-mail} e recuperação de senha.

\end{description}


\pretextualchapter{Glossário}

\begin{description}

\item[\textit{Antienumeration}] Padrão de segurança em que respostas a entradas inexistentes (\textit{e.g.}, \textit{e-mail} não cadastrado) são mantidas idênticas às de entradas válidas, evitando que um atacante descubra usuários cadastrados por meio de respostas distintas.

\item[\textit{API Key}] Cabeçalho HTTP (\texttt{x-api-key}) validado pelo \textit{API Gateway} antes da invocação da \textit{Lambda}, utilizado para aplicar \textit{throttling} e quota.

\item[\textit{AsyncStorage}] \textit{API} de armazenamento \textit{key-value} persistente do \textit{React Native}, utilizada no aplicativo para guardar o \textit{JWT}, o usuário autenticado e a lista de favoritos.

\item[\textit{Cold start}] Tempo adicional de inicialização sofrido por uma \textit{Lambda} quando não há contêiner aquecido disponível, antes do código de usuário começar a executar.

\item[Evento] Atividade pública (cultural, esportiva, gastronômica, religiosa ou comunitária) realizada em um dos seis municípios-alvo do projeto, identificada e estruturada pelo \textit{pipeline} de extração via \textit{LLM}.

\item[\textit{id\_token}] Credencial no formato \textit{JWT}, assinada em \textit{RS256}, emitida pelo \textit{Google} após autenticação \textit{OAuth} bem-sucedida; o \textit{backend} valida sua assinatura e troca-a pelo \textit{JWT} do próprio sistema.

\item[\textit{isEvent}] Atributo booleano definido pelo processador de linguagem natural para cada documento bruto: \texttt{true} indica que a página foi classificada como um evento e fica visível no aplicativo; \texttt{false} indica notícia comum, descartada pela \textit{API}.

\item[\textit{JWT}] Credencial \textit{stateless} no formato \textit{JSON Web Token}, assinada em \textit{HS256}, que prova a identidade do usuário ao \textit{backend} por sete dias após o \textit{login}.

\item[\textit{Linking} automático] Comportamento do \textit{login} com \textit{Google} em que, se o \textit{e-mail} do \textit{id\_token} já corresponde a uma conta existente, o identificador \textit{googleId} é vinculado a essa conta sem solicitar senha.

\item[\textit{Pipeline}] Sequência de componentes \textit{serverless} do Rolê SC que executa, de forma encadeada e idempotente, a coleta dos sites oficiais, o armazenamento bruto e a extração estruturada de eventos.

\item[\textit{Presigned URL}] \textit{URL} temporária do \textit{Amazon S3} (com tempo de vida de cinco minutos) assinada pelo \textit{backend}, que autoriza o aplicativo a executar um \textit{PUT} diretamente no \textit{bucket} sem que o conteúdo do arquivo passe pelo servidor.

\item[\textit{Raw crawl}] Documento bruto (\textit{HTML} ou texto) capturado de um \textit{site} público pelo \textit{web crawler} e armazenado na coleção \texttt{raw\_crawls} antes do processamento de linguagem natural.

\item[\textit{Serverless}] Modelo de execução em nuvem no qual o provedor gerencia integralmente a infraestrutura subjacente, e o consumidor é cobrado apenas pelo tempo efetivo de execução das funções, sem provisionar servidores.

\item[\textit{Strength meter}] Componente visual que apresenta ao usuário, em tempo real, quais requisitos de força da senha (comprimento, presença de caracteres especiais, dígitos etc.) já foram atendidos.

\item[\textit{TTL Index}] Índice especial do \textit{MongoDB} que remove automaticamente, em segundo plano, documentos cujo campo de expiração já passou, utilizado neste trabalho para tokens efêmeros de verificação de \textit{e-mail} e recuperação de senha.

\end{description}


\pretextualchapter{Glossário}

\begin{description}

\item[\textit{Antienumeration}] Padrão de segurança em que respostas a entradas inexistentes (\textit{e.g.}, \textit{e-mail} não cadastrado) são mantidas idênticas às de entradas válidas, evitando que um atacante descubra usuários cadastrados por meio de respostas distintas.

\item[\textit{API Key}] Cabeçalho HTTP (\texttt{x-api-key}) validado pelo \textit{API Gateway} antes da invocação da \textit{Lambda}, utilizado para aplicar \textit{throttling} e quota.

\item[\textit{AsyncStorage}] \textit{API} de armazenamento \textit{key-value} persistente do \textit{React Native}, utilizada no aplicativo para guardar o \textit{JWT}, o usuário autenticado e a lista de favoritos.

\item[\textit{Cold start}] Tempo adicional de inicialização sofrido por uma \textit{Lambda} quando não há contêiner aquecido disponível, antes do código de usuário começar a executar.

\item[Evento] Atividade pública (cultural, esportiva, gastronômica, religiosa ou comunitária) realizada em um dos seis municípios-alvo do projeto, identificada e estruturada pelo \textit{pipeline} de extração via \textit{LLM}.

\item[\textit{id\_token}] Credencial no formato \textit{JWT}, assinada em \textit{RS256}, emitida pelo \textit{Google} após autenticação \textit{OAuth} bem-sucedida; o \textit{backend} valida sua assinatura e troca-a pelo \textit{JWT} do próprio sistema.

\item[\textit{isEvent}] Atributo booleano definido pelo processador de linguagem natural para cada documento bruto: \texttt{true} indica que a página foi classificada como um evento e fica visível no aplicativo; \texttt{false} indica notícia comum, descartada pela \textit{API}.

\item[\textit{JWT}] Credencial \textit{stateless} no formato \textit{JSON Web Token}, assinada em \textit{HS256}, que prova a identidade do usuário ao \textit{backend} por sete dias após o \textit{login}.

\item[\textit{Linking} automático] Comportamento do \textit{login} com \textit{Google} em que, se o \textit{e-mail} do \textit{id\_token} já corresponde a uma conta existente, o identificador \textit{googleId} é vinculado a essa conta sem solicitar senha.

\item[\textit{Pipeline}] Sequência de componentes \textit{serverless} do Rolê SC que executa, de forma encadeada e idempotente, a coleta dos sites oficiais, o armazenamento bruto e a extração estruturada de eventos.

\item[\textit{Presigned URL}] \textit{URL} temporária do \textit{Amazon S3} (com tempo de vida de cinco minutos) assinada pelo \textit{backend}, que autoriza o aplicativo a executar um \textit{PUT} diretamente no \textit{bucket} sem que o conteúdo do arquivo passe pelo servidor.

\item[\textit{Raw crawl}] Documento bruto (\textit{HTML} ou texto) capturado de um \textit{site} público pelo \textit{web crawler} e armazenado na coleção \texttt{raw\_crawls} antes do processamento de linguagem natural.

\item[\textit{Serverless}] Modelo de execução em nuvem no qual o provedor gerencia integralmente a infraestrutura subjacente, e o consumidor é cobrado apenas pelo tempo efetivo de execução das funções, sem provisionar servidores.

\item[\textit{Strength meter}] Componente visual que apresenta ao usuário, em tempo real, quais requisitos de força da senha (comprimento, presença de caracteres especiais, dígitos etc.) já foram atendidos.

\item[\textit{TTL Index}] Índice especial do \textit{MongoDB} que remove automaticamente, em segundo plano, documentos cujo campo de expiração já passou, utilizado neste trabalho para tokens efêmeros de verificação de \textit{e-mail} e recuperação de senha.

\end{description}


\end{document}
